\chapter{Apriori: Assoziationsregeln}
\label{chap:apriori}

\begin{myremark}[Algorithmusart in diesem Kapitel]
Funktionsart: \textbf{Assoziation}. \\
Umsetzungsart: \textbf{datengetrieben}.
\end{myremark}

\begin{myremark}
Dieses Kapitel zeigt ein anderes Ziel als Klassifikation oder Regression:
Wir suchen nicht primär eine Vorhersage, sondern \textbf{häufige Muster} und \textbf{Regeln}
in Transaktionsdaten.
\end{myremark}

\section{Grundidee: Von häufigen Itemsets zu Regeln}

\begin{mydefinition}[Apriori]
Der \textbf{Apriori-Algorithmus} findet in Warenkörben (oder ähnlichen Datensätzen)
\textbf{häufige Itemsets} und daraus \textbf{Assoziationsregeln} wie
\enquote{Wenn A und B vorkommen, tritt oft auch C auf}.

Die Kernidee lautet: Wenn ein Itemset häufig ist, dann sind auch alle seine Teilmengen häufig.
Umgekehrt gilt: Ist eine Teilmenge nicht häufig, kann keine Obermenge häufig sein.
\end{mydefinition}

\section{Zentrale Kennzahlen}

Für eine Regel $X \Rightarrow Y$ gelten:
\[
\text{Support}(X \Rightarrow Y) = P(X \cup Y)
\]
\[
\text{Confidence}(X \Rightarrow Y) = \frac{P(X \cup Y)}{P(X)}
\]
\[
\text{Lift}(X \Rightarrow Y) = \frac{P(X \cup Y)}{P(X)\,P(Y)}
\]

\begin{itemize}
    \item \textbf{Support} misst, wie häufig die Kombination insgesamt ist.
    \item \textbf{Confidence} misst, wie zuverlässig die Regel ist.
    \item \textbf{Lift} misst, ob der Zusammenhang stärker ist als Zufall
          (Lift $> 1$ spricht für positive Assoziation).
\end{itemize}

\begin{myexercise}[Kennzahlen berechnen]
Ein Shop hat 200 Transaktionen:
\begin{itemize}
    \item In 50 Transaktionen sind \enquote{Pasta} und \enquote{Tomatensauce} gemeinsam enthalten.
    \item \enquote{Pasta} kommt in 80 Transaktionen vor.
    \item \enquote{Tomatensauce} kommt in 70 Transaktionen vor.
\end{itemize}
Berechnen Sie für die Regel
\enquote{Pasta $\Rightarrow$ Tomatensauce}:
\begin{enumerate}
    \item Support
    \item Confidence
    \item Lift
\end{enumerate}
\begin{myanswer}
Support $= \frac{50}{200}=0.25$. \\
Confidence $= \frac{50}{80}=0.625$. \\
Lift $= \frac{0.25}{0.4\cdot0.35}=1.79$ (deutlich über 1).
\end{myanswer}
\end{myexercise}

\section{Warum Schwellenwerte wichtig sind}

\begin{mycaution}
Ohne sinnvolle Grenzwerte (z.\,B. \texttt{min\_support}, \texttt{min\_confidence})
produziert Apriori sehr viele triviale oder zufällige Regeln.
Das führt schnell zu \textbf{Scheinkorrelationen}.
\end{mycaution}

\begin{myexercise}[Regelqualität bewerten]
Ein Data-Team meldet die Regel:
\enquote{Energy-Drink $\Rightarrow$ Kopfhörer},
Confidence $= 0.82$, aber Support $= 0.01$.
\begin{enumerate}
    \item Warum ist diese Regel trotz hoher Confidence kritisch zu prüfen?
    \item Welche zusätzlichen Prüfungen würden Sie durchführen?
\end{enumerate}
\end{myexercise}

\section{Anwendungsbezug}

\begin{mychallenge}[Apriori in der Praxis]
Skizzieren Sie ein Mini-Projekt für Apriori mit einem Datensatz aus Schule oder Alltag.
Beantworten Sie:
\begin{enumerate}
    \item Was ist eine Transaktion?
    \item Was sind die Items?
    \item Welche Entscheidung soll mit den Regeln unterstützt werden?
    \item Welche Risiken von Fehlinterpretationen sehen Sie?
\end{enumerate}
\end{mychallenge}
